% =====================================================================
% REFRAME DRAFT (narrow + honest) — staging only, NOT \input by main.
% Numbers in ⟦TOKENS⟧ are re-run-dependent; populate from outputs/duke_sweeps_fixed
% once the corrected sweep lands.  Legend at the bottom gives each token's
% CURRENT (pre-correction) value so you can sanity-check the drop-in.
% Addresses reviewer Major #2 (within-species / fraction), #5 (optimize-then-
% validate), and the #7-driven geometric-ceiling rewording.
% =====================================================================

% ---------------------------------------------------------------------
% (A) ABSTRACT — replace the deployment-penalty sentences.
% ---------------------------------------------------------------------
Using FEM-derived fields over a 12-contact cuff on 29 histology-segmented
vagus nerves (18 swine, 11 human), we find that configurations optimized on
sparsely-sampled fibers and \emph{deployed without re-scoring on the full
population} lose substantial selectivity --- a \emph{deployment penalty}
invisible to reduced-order models.  The penalty is small in swine
(median ⟦PEN_SW_MED⟧) but large in human (⟦PEN_HU_MED⟧; Mann--Whitney
⟦P_SPECIES⟧), tracking the number of fibers per target fascicle
(Spearman $r=$⟦RHO_POOL⟧ across nerves), which differs systematically
between species; it manifests as failure to recruit the large human target
fascicle (⟦P_RECRUIT⟧) rather than off-target spill (n.s.).  Crucially, the
loss is a cost of \emph{not} validating on the population: re-scoring the
sparse-optimized design on the full population --- which our simulator makes
cheap --- recovers most of it (residual human penalty ⟦PEN_HU_OTV⟧), so the
practical prescription is to optimize and/or validate at population scale,
now feasible.  Reduced-order fascicular sampling therefore overstates the
selectivity that will actually be delivered when used the standard way, most
severely on the human vagus geometry most relevant to clinical translation.

% ---------------------------------------------------------------------
% (B) RESULTS §3.5 — reframed sparse-sampling subsection.
% ---------------------------------------------------------------------
\subsection{Reduced-order fiber sampling overestimates deliverable
  selectivity unless the design is re-scored on the population}
\label{sec:results-sparse}

Prior GPU-accelerated selectivity work reduces each fascicle to a single
centroid fiber to keep the population tractable~\cite{hussain2024}: a design
is optimized on the reduced set and the reported selectivity is read off that
same reduced set, without re-scoring on the fibers that are not simulated.
Because jaxon simulates the full ${\sim}1000$-fiber population directly, we
can ask what that approximation costs at deployment.  For each nerve $\times$
seed we re-ran the identical optimization on four sparse fiber-sampling levels
(one centroid fiber per fascicle, and $1$, $3$ or $10$ randomly-sampled fibers
per fascicle) by pure row-subsampling of the precomputed $V_e$ lead-field, and
re-scored each sparse-optimized amplitude vector on the full population
\emph{without re-optimizing}, giving the transfer selectivity
$\mathrm{SI_{transfer}}$ --- the selectivity actually delivered.  The
deployment penalty is the drop from the dense ceiling,
$\mathrm{SI_{dense}} - \mathrm{SI_{transfer}}$
(Figure~\ref{fig:duke}\textbf{d}).

\paragraph{The penalty is small in swine, large in human.}
Per nerve it is a median ⟦PEN_SW_MED⟧ in swine (mean ⟦PEN_SW_MEAN⟧) versus
⟦PEN_HU_MED⟧ in human (mean ⟦PEN_HU_MEAN⟧; Mann--Whitney ⟦P_SPECIES⟧), and the
transfer SI lies significantly below the dense ceiling for every sampling
strategy in both species (two-sided per-nerve Wilcoxon, Holm-corrected; all
⟦P_TRANSFER⟧).

\paragraph{Adding fibers does not close the gap; sampling \emph{fraction}
  does.}
Sampling \emph{density} in absolute counts barely matters: ten fibers per
fascicle transfer no better than one, or than the centroid.  Re-expressed as
the fraction of each target fascicle that is sampled --- the species-fair axis,
since ten fibers is ${\sim}\!⟦FRAC_SW⟧$ of a typical swine fascicle but only
${\sim}\!⟦FRAC_HU⟧$ of a human one --- the penalty is concentrated below
${\sim}⟦FRAC_KNEE⟧$ of the fascicle population and is essentially closed above
it; the species gap persists at matched low fraction (human ⟦PEN_HU_LOWFRAC⟧
vs swine ⟦PEN_SW_LOWFRAC⟧), so the effect is not merely a sampling-fraction
artifact: a large human fascicle needs a larger \emph{absolute} number of
sampled fibers to be represented.

\paragraph{Fascicle size sets the between-species difference.}
The per-nerve penalty correlates with the number of fibers per target fascicle
across the cohort (Spearman $r=$⟦RHO_POOL⟧, ⟦P_POOL⟧, $n=29$ nerves), which has
median ⟦FPF_SW⟧ in swine versus ⟦FPF_HU⟧ in human.  This anatomical axis
explains the \emph{between-species} difference; within each species the
correlation is directionally consistent but not individually significant at
these sample sizes (swine $r=$⟦RHO_SW⟧, ⟦P_RHO_SW⟧, $n=18$; human
$r=$⟦RHO_HU⟧, ⟦P_RHO_HU⟧, $n=11$), so we frame fibers-per-fascicle as the
driver of the species gap rather than as a within-species predictor.  Swine
vagus nerves carry many small fascicles, so a handful of sampled fibers
represents the target well and the optimized field transfers; human nerves
carry a few large fascicles, so a field tuned to one sampled fiber does not
cover the rest of that fascicle.  Counterintuitively, the more fascicular
nerve (swine) is the more forgiving one.

\paragraph{The mechanism is failure to recruit, not off-target spill.}
Relative to the dense optimum, the sparse-transfer configuration in human loses
on-target recruitment (median ⟦REC_HU_DENSE⟧ $\rightarrow$ ⟦REC_HU_TRANS⟧;
⟦P_RECRUIT⟧) and selectivity (⟦P_SI_LOSS⟧), whereas off-target recruitment and
total delivered current are statistically unchanged (⟦P_OFF⟧)
(Figure~\ref{fig:duke}\textbf{b}).

\paragraph{Re-scoring on the population recovers most of the penalty.}
The penalty above is the cost of the standard workflow --- optimize on few
fibers, deploy without re-scoring.  Because jaxon can re-score a candidate on
the full population cheaply, we asked what remains if one instead selects, per
nerve, the sparse-optimized design that scores best on the full population
(optimize-then-validate).  The penalty then largely collapses: swine
⟦PEN_SW_MED⟧ $\rightarrow$ ⟦PEN_SW_OTV⟧ and human ⟦PEN_HU_MED⟧ $\rightarrow$
⟦PEN_HU_OTV⟧.  The honest conclusion is therefore narrower than ``reduced-order
models are wrong'': centroid-only sampling \emph{as standardly used} ---
optimized and reported on the reduced set~\cite{hussain2024} --- materially
overstates deliverable selectivity on large-fascicle (human) anatomies, but a
practitioner who re-scores or selects on the full population, which is now
computationally feasible, recovers most of it.  A residual species-dependent
penalty (${\sim}⟦PEN_HU_OTV⟧$ in human) survives even optimize-then-validate,
because no sparse design fully covers a large fascicle.  The practical
prescription is to bring the full population into the loop --- at optimization
or at least at validation --- precisely the regime jaxon's throughput opens up.

% ---------------------------------------------------------------------
% (C) RESULTS §3.4 — reworded "Geometric-ceiling cases" paragraph.
% (#7: euclidean fascicle–contact distance is only a weak SI predictor, ρ≈−0.13
%  n.s.; the hard cases are large-fascicle, not equidistant-fascicle.)
% ---------------------------------------------------------------------
\paragraph{Geometric-ceiling cases.}
A minority of human anatomies cannot be steered well even on the full
population.  These are not primarily cases where off-target fascicles happen to
sit as close to the contacts as the target --- across the cohort the minimum
fascicle-to-contact distance is only a weak predictor of achieved selectivity
(Spearman ⟦RHO_DIST⟧, n.s.) --- but cases where the target fascicle is large
relative to the steerable field, so that sparing every off-target fiber while
recruiting the whole target is geometrically impossible.  These
geometric-ceiling cases form the heavy lower tail of the human distribution and
are \emph{retained} in all statistics --- they are genuine selectivity outcomes
on selectable targets, not solver artifacts.  Separately, we exclude five seeds
(all human) whose target cluster spans more than $60\,\%$ of the nerve
cross-section (target fraction $> 0.6$): such targets are degenerate by
construction --- no selective solution exists --- so they are uninformative
about achievable selectivity.

% ---------------------------------------------------------------------
% (D) DISCUSSION — reworded opening deployment-penalty framing.
% ---------------------------------------------------------------------
Using population-scale differentiable simulation, we have shown that the
reduced-order fiber sampling standard in GPU-accelerated selectivity modeling
--- reducing each fascicle to one or a few representative fibers and reading
selectivity off that reduced set --- overstates the selectivity actually
delivered to the nerve when used that way.  The error is not a fixed offset:
this \emph{deployment penalty} is governed by fascicular architecture --- the
number of fibers per target fascicle --- and is therefore small in swine
(median ⟦PEN_SW_MED⟧) but large in human (median ⟦PEN_HU_MED⟧), a
between-species translational bias that reduced-order models cannot detect
because they never simulate the fibers the bias acts on.  Mechanistically the
penalty is a failure to recruit the large human target fascicle from a field
tuned to a few sampled fibers, not off-target spill or reduced drive.  It is
also largely \emph{recoverable}: because the full-population forward pass is
cheap in jaxon, re-scoring or selecting a sparse-optimized design on the
population removes most of the penalty (residual human ${\sim}⟦PEN_HU_OTV⟧$),
so the actionable message is that population-scale validation --- not a
different optimizer --- is what reduced-order pipelines lack.  This was
measurable only because jaxon simulates the full ${\sim}1000$-fiber population
\dots

% =====================================================================
% PLACEHOLDER LEGEND — token : current (pre-correction) value / source.
% Repopulate ALL from outputs/duke_sweeps_fixed via the reviewer re-analysis.
% =====================================================================
% ⟦PEN_SW_MED⟧   swine penalty median        current 0.03  (likely ~unchanged)
% ⟦PEN_SW_MEAN⟧  swine penalty mean          current 0.04
% ⟦PEN_HU_MED⟧   human penalty median        current 0.17  (EXPECTED TO DROP)
% ⟦PEN_HU_MEAN⟧  human penalty mean          current 0.23  (EXPECTED TO DROP)
% ⟦PEN_SW_OTV⟧   swine optimize-then-validate current 0.001
% ⟦PEN_HU_OTV⟧   human optimize-then-validate current 0.067
% ⟦P_SPECIES⟧    species MWU p               current p<0.001
% ⟦P_TRANSFER⟧   all-cells Wilcoxon          current p<=0.005
% ⟦RHO_POOL⟧     pooled Spearman             current +0.73
% ⟦P_POOL⟧       pooled p                    current p<0.001
% ⟦RHO_SW⟧/⟦P_RHO_SW⟧  within swine          current +0.28 / p=0.26
% ⟦RHO_HU⟧/⟦P_RHO_HU⟧  within human          current +0.44 / p=0.18
% ⟦FPF_SW⟧/⟦FPF_HU⟧    fibers/target fasc median  current 30 / 151
% ⟦FRAC_SW⟧/⟦FRAC_HU⟧  10-fiber fraction      current ~33% / ~7%
% ⟦FRAC_KNEE⟧    fraction where penalty closes current ~2%
% ⟦PEN_HU_LOWFRAC⟧/⟦PEN_SW_LOWFRAC⟧  penalty at matched low fraction  current ~0.10 / ~0.00
% ⟦REC_HU_DENSE⟧/⟦REC_HU_TRANS⟧  human on-target recruit  current 100% / 65%
% ⟦P_RECRUIT⟧    recruit-loss p             current p=0.012 (abstract uses 0.01)
% ⟦P_SI_LOSS⟧    SI-loss p                  current p=0.007
% ⟦P_OFF⟧        off-target/current p       current p>0.7
% ⟦RHO_DIST⟧     SI vs min fasc-contact dist Spearman  current -0.13 (n.s., n=111)
